% FILE: 01_research_question.tex
% PROJECT: phd-thesis
% MANUFACTURED: 2026-06-01
% AUTHORITY: phd-thesis/ledgers/PROJECT_INDEX.yaml, checkpoints/PI_RESEARCH_PROGRAM_ALIVE_001.md
% DO NOT CLAIM: See ledgers/DO_NOT_CLAIM_LEDGER.md

\section{Research Question}
\label{sec:phd-thesis-question}

\subsection{Problem Statement}
\label{sec:phd-thesis-question-problem}

The central problem this dissertation addresses is the systematic gap between what enterprise
software systems \emph{declare} they do and what the event log \emph{proves} they do. This gap
is not a logging deficiency or a monitoring failure; it is a structural property of systems that
optimize for prediction and automation while leaving process consequence --- the lawful ordering
of activities, the documented transfer of custody, the admissible evidence of conformance ---
unmeasured and therefore unmanageable.

The problem manifests across three observable failure modes. First, conformance theater: a
system reports process compliance without a named witness, without a fitness score derived from
replay against a declared model, and without a receipt binding the claim to a specific event log.
Second, metric fabrication: EBITDA projections, rework rates, and throughput figures are
presented to boards without the process model hash and log hash that would allow independent
verification. Third, variant explosion: the declared process is a single trace while the
discovered process --- mined from the actual event log --- reveals hidden loops, rework paths,
and unauthorized shortcuts that were invisible to the system's own reporting layer.

The Van der Aalst Constitution (\texttt{\textasciitilde/.claude/rules/process-mining-chicago-tdd.md})
states this as a hostile assumption: ``The declared manufacturing pipeline is not the real
runtime process.'' The dissertation takes this assumption as its foundational axiom and asks
what infrastructure is required to close the gap.

\subsection{Old-Regime Assumption Challenged}
\label{sec:phd-thesis-question-old-regime}

The old-regime assumption under challenge is that prediction is a sufficient substitute for
coordination. In the prediction regime, an AI or rule-based system forecasts the next state of
a process and acts on that forecast. The process model is implicit in the system's weights or
business rules; the event log is a side channel, not the primary authority. Conformance is
inferred from accuracy metrics rather than derived from replay.

This assumption fails when the consequences of process execution become legally, financially,
or operationally binding. A board receiving an M\&A diligence package cannot accept a
conformance claim grounded in prediction accuracy; it requires a named witness, a process model
hash, an event log hash, and a fitness score computed by replaying the log against the model.
A regulatory audit cannot accept an SLA compliance claim grounded in a monitoring dashboard;
it requires an object-centric event log traceable to individual process objects and activities.

The dissertation documents this failure through the 2016 language-model experiment (see
Section~\ref{sec:phd-thesis-lineage-2016}), which demonstrated that local text imitation does
not conserve consequence --- a finding that motivated the entire research program recorded in
this foundry. The thesis lineage, as captured in \texttt{ledgers/PROJECT\_INDEX.yaml}, traces
directly from that experiment through the Chatman Equation, receipt infrastructure, ggen,
wasm4pm, and the full-lifecycle process intelligence architecture.

\subsection{Hypothesis}
\label{sec:phd-thesis-question-hypothesis}

The hypothesis of this dissertation is: \emph{Full-lifecycle process intelligence --- defined
as the capacity to ingest, validate, replay, discover, conform, simulate, repair, optimize,
monitor, and decommission process evidence under a single, receipted type-law --- is a
necessary and sufficient foundation for board-admissible process claims.}

This hypothesis is operationalized through the Chatman Equation
$A = \mu(O^*)$, which asserts that process authority ($A$) is the application of the process
mining function ($\mu$) to the full object-centric process state ($O^*$). The equation is
cited in the thesis lineage recorded in \texttt{ledgers/PROJECT\_INDEX.yaml} and in the
mathematical invariants listed in checkpoint
\texttt{PROCESS\_INTELLIGENCE\_ALIVE\_001.md}, which includes the admissibility boundary
$R \vdash P_i = \mu(O^*, T, L)$.

The hypothesis implies three testable claims: (1)~every process claim can be grounded in a
specific event log, process model, and fitness score; (2)~every receipt chain is
cryptographically verifiable and can be audited independently of the system that emitted it;
(3)~every gap in the process evidence --- missing activities, unauthorized variants, object
lifecycle violations --- can be identified, documented, and traced to a structural defect
rather than a data quality problem.

\subsection{Success Criteria}
\label{sec:phd-thesis-question-success}

Success for this dissertation is defined by the ALIVE verdict criteria recorded in
\texttt{phd-thesis/README.md}. Specifically, the dissertation artifact is ALIVE when:

\begin{enumerate}
  \item TeX source files exist in the declared locations under \texttt{chapters/} and
    \texttt{frontmatter/}.
  \item The source ledger (\texttt{ledgers/PROJECT\_INDEX.yaml}) is present and populated ---
    a condition already met (34 projects catalogued as of 2026-06-01).
  \item The claim ledger prohibitions (\texttt{ledgers/DO\_NOT\_CLAIM\_LEDGER.md}) have been
    respected throughout manufacture --- verified by the absence of prohibited terms in all
    artifacts.
  \item The PDF compiles without error from the TeX sources using pdfTeX 3.141592653-2.6-1.40.26
    (TeX Live 2024), as confirmed available in \texttt{ledgers/WORKFLOW\_RECEIPT.yaml}.
  \item The compiled PDF hash is recorded in the receipt chain, binding the final artifact to
    the specific source commit and toolchain invocation that produced it.
\end{enumerate}

The broader research program has met its own ALIVE criteria: PI\_RESEARCH\_PROGRAM\_ALIVE\_001
sealed 2026-06-01 with 12/12 audit gates PASS across 342 research artifacts. The dissertation
workspace is PARTIAL pending TeX manufacture and PDF compilation, which are the next receipted
work items.
